% Compile with: xelatex -output-directory=report report/main.tex
\documentclass[UTF8,a4paper,11pt]{ctexart}
\usepackage{amsmath}
\usepackage{booktabs}
\usepackage{geometry}
\usepackage{graphicx}
\usepackage{hyperref}
\geometry{margin=2.5cm}
\hypersetup{colorlinks=true,linkcolor=black,urlcolor=blue,citecolor=black}

\title{基于 LSTM 与 GRU 的唐诗字符级语言建模与生成研究}
\author{Houkiang}
\date{日期待填}

\begin{document}
\maketitle

\begin{abstract}
\noindent\textbf{Abstract.}
研究问题、方法、数据规模、测试集困惑度和主要发现待实验完成后填写。
\end{abstract}

\noindent\textbf{Keywords:} 字符级语言模型；LSTM；GRU；困惑度；文本生成

\section{Introduction}
说明字符级唐诗建模的研究问题、比较目标与本文贡献。数据与任务背景见\cite{tangdata}。

\section{Related Work}
讨论循环神经网络、LSTM、GRU 和序列生成评价的相关工作；教材依据见\cite{nndl2020}。补充原始论文引用与准确定位。

\section{Method}
定义字符词表、序列边界、模型结构、优化目标与采样策略。主要评价量为字符级困惑度：
\begin{equation}
  \mathrm{PPL}=\exp\!\left(\frac{1}{N}\sum_{i=1}^{N}-\log p_{\theta}(y_i\mid y_{<i})\right).
\end{equation}
其中 $N$ 只计被预测字符与诗篇结束标记，不计起始标记和填充位。

\section{Experiment}
记录语料 SHA-256、获取日期、划分篇数、词表与 OOV、超参数、训练预算、硬件、软件版本、随机种子及基线。主实验应保持 LSTM 与 GRU 的数据协议一致。

\section{Result}
实验完成后插入测试集交叉熵与困惑度表、训练曲线、多个固定种子的生成样本、约束解码对照和必要消融。每个图表附对应 run ID；不展示未验证数值。

\section{Conclusion}
根据真实结果总结模型比较、生成特征、数据与评估局限，以及可复现的后续工作。

\begin{thebibliography}{9}
\bibitem{nndl2020}
邱锡鹏. \textit{神经网络与深度学习}. 机械工业出版社, 2020. ISBN 978-7-111-64968-7.
\bibitem{tangdata}
nndl. \textit{poetryFromTang.txt}. \url{https://github.com/nndl/llm-beginner/blob/master/poetryFromTang.txt}. 访问日期待填.
\end{thebibliography}

\end{document}
